\documentclass[12pt]{article}
\usepackage[utf8]{inputenc}
\usepackage[T1]{fontenc}
\usepackage{hyperref}
\title{Tapesena Language Resources for K-12 Curriculum Development}
\author{Compiled by Claude Sonnet 5 (claude-sonnet-5) for Andr\'as Kornai}
\date{}
\begin{document}
\maketitle

\section{Introduction}
Tapesena (ISO 639-3 vnp; Glottocode vuna1239) is an Oceanic language of northern Espiritu Santo, Vanuatu. In the sources found it is listed under the name Vunapu, with Tapesena given as an alternative name and Tapechalil as a dialect name. Speaker figures (about 380) date from 1983, and sources disagree on its vitality. Very little exists for it, and no school materials were found.

\section{Primary Educational Resources}
No primers, textbooks or school programmes for Tapesena were found in search results.

\section{Linguistic and Reference Materials}
General reference entries only:
\begin{itemize}
\item Wikipedia, Vunapu language: \url{https://en.wikipedia.org/wiki/Vunapu_language}
\item Glottolog entry (vuna1239): \url{https://glottolog.org/resource/languoid/id/vuna1239.bigmap.html}
\item Verbix Endangered Languages entry: \url{https://docs.verbix.com/EndangeredLanguages/Vunapu}
\end{itemize}

\section{Religious and Community Resources}
Joshua Project lists the people group and notes that a Bible translation is still needed:
\url{https://joshuaproject.net/people_groups/15805}, \url{https://joshuaproject.net/languages/vnp}

\section{Digital and Audio Resources}
None found.

\section{Curriculum Development Context}
No documentation suitable for curriculum use was found. Any curriculum would need new work with the community, starting from fieldwork-based description.
\end{document}